\section{The conditional fold that pays for the new source node}

The recompiled source using the $C$ parameter in the generalization row has
202 distinct source subtrees. A corrected literal allocation first produces
216 working states. The required return to 215 states uses a new conditional
decoder fold; no existing conditional-marker helper is copied into a spare
state. An already present decoder chain supplies the required dispatch.

Let $T$ denote the corrected 216-state table and $N$ the resulting 215-state
table. Define $\phi(310)=129$ and $\phi(q)=q$ otherwise. The construction deletes
the two rows of state 310, redirects every target 310 to 129, and replaces
the remaining row $129/1$ by $1L216$. All other rows are retained, subject
only to the target redirection. In this table the sole redirected incoming
row is $312/0$.

\begin{center}
\begin{tabular}{c|cc}
Original state & Read 0 & Read 1\\\hline
129 & $0R154$ & $1R189$\\
310 & $0R154$ & $1R86$
\end{tabular}
\end{center}

The complete incoming-row census is particularly simple. State 129 has only
the incoming row $211/1=1R129$, and state 310 has only the incoming row
$312/0=0R310$. Neither state is the initial control. Therefore, on every
execution starting in control 0, an arrival at old state 129 has a one
immediately to its left, while an arrival at old state 310 has a zero there.
This statement holds on arbitrary initial binary tapes.

Both old roles agree on read 0. On read 1, the new row moves left to state
216 and reads this origin marker. For marker one, its execution is
\[
 129 \xrightarrow{1L}216\xrightarrow{1R}178\xrightarrow{1R}189.
\]
For marker zero, the execution is
\[
 129\xrightarrow{1L}216\xrightarrow{0R}157
     \xrightarrow{1L}148\xrightarrow{0R}186\xrightarrow{1R}86.
\]
Every write reproduces the symbol already present. Both paths end one cell
to the right of the original head, in exactly the control required by the
corresponding old role. No exterior tape cell is read. Thus one old step
is simulated by respectively three or five new steps on read 1, and by one
new step on read 0.

At every other old state, one step commutes with $\phi$. Induction therefore
matches the complete old computation from control 0 with the new computation
at a sequence of boundaries having identical tapes and heads and controls
related by $\phi$. Each matching block has positive length at most five.
Consequently $T$ and $N$ halt together, on every initial tape from control 0.
The claim is not equivalence from arbitrary old control 310 with an arbitrary
unconstrained left marker.

The independent checker reads the literal tables, verifies the complete
incoming census and all four local blocks, checks all redirected rows, and
checks that all 215 surviving working states remain graph reachable. It also
produces a dense table with working controls $0,\ldots,214$ and separate halt
control 215, and verifies the exact state-label bijection. The native table
uses separate halt control 254. Universal primitive simulation and the blank
bootstrap are independently checked obligations; the finite fold proof alone
does not establish their premises.

One allocator repair is essential to this construction. Row $134/0=0R17$
must be protected because the existing engine can enter it through
$11/0=0L26$ and $26/0=0R134$. A prefix-only decoder audit does not detect a
change to that engine use. The uncorrected 216-state allocation failed its
first arbitrary-natural primitive contract, and is not a candidate.
